\centerline{\bf \Huge Степень точки}
\centerline{\bf \Huge Level: Hard}

\begin{flushright}
        \scriptsize{Ты --- это ты. Невелика разница между тем,\\
                    как ты воспринимаешь себя и как тебя воспринимают другие.\\
                    Рей Аянами.}
\end{flushright}

\problem{1.1}
В окружности $\Omega$ проведена хорда $BC$. В один из отсекаемых ею сегментов вписаны две окружности $\omega_1$ и $\omega_2$. Докажите, что середина дуги $BC$ окружности $\Omega$, не касающейся окружностей $\omega_1$ и $\omega_2$, лежит на радикальной оси окружностей $\omega_1$ и $\omega_2$.

\problem{1.2}Окружности $\Omega$ и $\omega$ касаются друг друга внутренним образом в точке $A$. Проведем в большей окружности $\Omega$ хорду $BC$, касающуюся $\omega$ в точке $K$ (хорда $BC$ не является диаметром $\Omega$). Точка $M$ — середина отрезка $AK$. Докажите, что окружность $(BMC)$ проходит через центр $\omega$.

\problem{2}
В треугольнике $ABC$ проведены чевианы $AA_1$, $BB_1$ и $CC_1$, пересекающиеся в одной точке. Прямая $B_1C_1$ пересекает окружность $(ABC)$ в точках $P$ и $Q$. Докажите, что окружность $(A_1PQ)$ проходит через середину отрезка $BC$.

\problem{3}
В треугольнике $ABC$ отмечена точка пересечения высот $H$. Окружность с центром в середине стороны $AB$ и проходящая через $H$, пересекает сторону $AB$ в точках $C_1$ и $C_2$. Аналогично определяются точки $A_1$, $A_2$ и $B_1$, $B_2$. Докажите, что точки $A_1$, $A_2$, $B_1$, $B_2$, $C_1$, $C_2$ лежат на одной окружности.

\problem{4}
На стороне $AB$ выпуклого четырёхугольника $ABCD$ взяты точки $K$ и $L$ (точка $K$ лежит между $A$ и $L$), а на стороне $CD$ взяты точки $M$ и $N$ (точка $M$ между $C$ и $N$). Известно, что $AK=KN=ND$ и $LB=BC=CM$. Докажите, что если $BCNK$ — вписанный четырехугольник, то и $ADML$ тоже вписан.

\problem{5}
В остроугольном треугольнике $ABC$ высоты $AA_1$, $BB_1$ и $CC_1$ пересекаются в точке $H$. Точка $O$ — центр окружности $(ABC)$. Докажите, что точки $O$, $H$, $A_1$ и точка, симметричная точке $A$ относительно прямой $B_1C_1$, лежат на одной окружности.

\problem{6}
Окружность $\omega$ касается сторон угла $\angle BAC$ в точках $B$ и $C$. Прямая $\ell$ пересекает отрезки $AB$ и $AC$ в точках $K$ и $L$ соответственно. Окружность $\omega$ пересекает прямую $\ell$ в точках $P$ и $Q$. Точки $S$ и $T$ выбраны на отрезке $BC$ так, что $KS$ параллельно $AC$ и $LT$ параллельно $AB$. Докажите, что точки $P$, $Q$, $S$ и $T$ лежат на одной окружности.


\problem{7}
Для окружности $\omega$ и точки $X$ обозначим через $Pow_\omega X$ степень этой точки относительно окружности $\omega$. Докажите, что для любых прямых $AB$ и $CD$, пересекающихся в точке $P$, и любой точки $X$ имеет место равенство $$Pow_{(PAB)}(X)+Pow_{(PCD)}(X)=Pow_{(PAD)}(X)+Pow_{(PBC)}(X).$$

\problem{8}
Пусть $A$, $V_1$, $V_2$, $B$, $U_2$, $U_1$ — точки на окружности $\Omega$, выбранные в таком порядке, что $BU_2>AU_1>BV_2>AV_1$. Точка $X$ — произвольная точка на дуге $V_1V_2$, не содержащей точки $A$ и $B$. Пусть $XA\cap U_1V_1=C$, $XB\cap U_2V_2=D$. Докажите, что существует такая точка $K$, не зависящая от точки $X$, что степень этой точки относительно окружности $(XCD)$ постоянна при любом выборе точки $X$.